\documentclass[11pt]{article}
\usepackage{amsmath,amsthm,hyperref}
\pdfinfoomitdate=1
\pdftrailerid{}
\ifdefined\pdfsuppressptexinfo\pdfsuppressptexinfo=15\fi
\setlength{\textwidth}{6.15in}
\setlength{\oddsidemargin}{0.175in}
\setlength{\evensidemargin}{0.175in}
\setlength{\textheight}{8.7in}
\setlength{\topmargin}{-0.3in}
\providecommand{\SourceCommit}{SOURCE-NOT-BUILT}
\newcommand{\R}{\mathbf R}
\newcommand{\C}{\mathbf C}
\newcommand{\Z}{\mathbf Z}
\newcommand{\norm}[1]{\lVert#1\rVert}
\newcommand{\ip}[2]{\langle#1,#2\rangle}
\newcommand{\e}{\mathrm e}
\DeclareMathOperator{\tr}{tr}
\DeclareMathOperator{\im}{im}
\DeclareMathOperator{\ran}{ran}
\DeclareMathOperator{\Rea}{Re}
\newtheorem{theorem}{Theorem}[section]
\newtheorem{lemma}[theorem]{Lemma}
\newtheorem{proposition}[theorem]{Proposition}
\newtheorem{corollary}[theorem]{Corollary}
\theoremstyle{definition}
\newtheorem{definition}[theorem]{Definition}
\theoremstyle{remark}
\newtheorem{remark}[theorem]{Remark}

\title{Sharp Logarithmic Loss for Endpoint Observations\\
in a Discrete Coexact Model}
\author{Lluis Eriksson}
\date{12 November 2026\\\small Working paper draft}

\begin{document}
\maketitle
\begin{abstract}
We determine the asymptotic order of the minimum energy invisible to a fixed
canonical partition in a finite coexact vector-field model. On an even
three-dimensional torus with counting norm and checkerboard source weights
$H=cn^2$, $a=c/48$, the full endpoint cell observations leave a nonzero
kernel whose minimum Rayleigh quotient has order $c/\log n$. Explicit finite
upper and lower bounds are proved. The upper witness is real, lies in the
actual image of the forward curl adjoint, and annihilates every original
cell mean. It follows that finite penalties on these observations cannot
restore a size-uniform target $c$. The proof combines the complete restricted
Fourier spectrum, an original-row range theorem, finite Green estimates,
optimization over complementary singleton coefficients, and frequency-band
removal localized to the lowest nonzero component sector. We also give an
exact conditional-cost identity for the remaining rows of moving partitions.
The best canonical partition, including a quadratic moving cut, remains
unresolved. All arguments are finite and analytical. Green-function
logarithms and Schur minimization are classical; external priority for the
precise model-specific theorem is unestablished. No physical Yang--Mills
gap or continuum conclusion is asserted.
\end{abstract}

\section{Main result and relation to earlier work}
Our main conclusion is a two-sided $c/\log n$ estimate for the full
observation kernel of one admissible fixed partition. The distinction
between a fixed partition and the best partition is essential. A weak
lower bound alone would not refute a uniform gap; here a real invisible
field supplies the matching upper order. Conversely, failure of this
fixed partition does not establish failure for moving partitions.

This manuscript is a proposed linked revision of the programme working
paper \emph{Spectral Reduction and Logarithmic Observation Response in a
Discrete Coexact Model}, AIRR version~1~\cite{Eriksson284}. That version
established the spectrum, range, and direct-response bounds but explicitly
left optimized cancellation open. The present revision retains their
proofs and adds the full endpoint-kernel theorem and conditional residual
cost. The earlier version remains a distinct immutable artifact; its
assessments are not assessments of this revision.

Checkerboard weights couple two Fourier frequencies, while the actual
curl-adjoint image imposes two transverse polarizations and removes the
harmonic zero mode. Neither a scalar block alone nor injectivity on low
modes decides high-mode cancellation. Sections~\ref{sec:spectrum}--\ref{sec:cost}
retain all those restrictions and the original rows. The response estimates
in Sections~\ref{sec:whole}--\ref{sec:useful} supply finite constants.
Section~\ref{sec:endpoint} then handles complementary coefficients and
constructs a full-cell witness; Section~\ref{sec:residual} isolates the
remaining problem for moving cuts.

Two-dimensional logarithmic Green functions are classical. Kenyon's
Theorem~7.3 treats critical planar graphs~\cite{Kenyon}; Lawler and Limic
develop continuous-time walks and the two-dimensional potential kernel,
including Theorem~4.4.4~\cite{LawlerLimic}. Their graph, volume, and operator
normalizations are not identified with our finite coexact compression.
Schur minimization is also classical; the primary abstract of Corach,
Maestripieri, and Stojanoff explicitly relates shorted operators and
projections~\cite{Corach}. Every finite comparison used here is proved
below. These bounded primary-source comparisons do not establish priority
for the precise endpoint theorem; neither the logarithmic mechanism nor
the general minimization algebra is claimed as new.

The earlier programme working paper \emph{Coverage and coexactness in
checkerboard slice observations}~\cite{Eriksson274} concerns coverage,
projection issues and an ambient-field lower bound. Those results are not
presented as new here. Remark~\ref{rem:comparison} distinguishes its domain
from the present actual-coexact lower bound.
\section{The original model and observations}\label{sec:model}
Let $n\ge4$ be even and $X=(\Z/n\Z)^3$. All physical norms and adjoints use
counting measure. For $u:X\to\R^3$, put
\[
 \Delta_j u(x)=u(x+e_j)-u(x),\qquad
 (d_1u)_{j\ell}=\Delta_j u_\ell-\Delta_\ell u_j\quad(j<\ell).
\]
The norm on two-forms counts each pair $j<\ell$ once. Our actual domain is
$U=\im d_1^*$, not an abstract divergence-free space with freely added
constants. We prove its exact Fourier description in Section~\ref{sec:spectrum}.

There are two original channels, indexed by $\varepsilon\in\{-1,+1\}$.
Write $\chi(x)=(-1)^{x_1+x_2+x_3}$ and, initially, let $H\ge a>0$.
Set $A_0=(H+a)/2$, $B_0=(H-a)/2$, and
\begin{equation}\label{eq:originalQ}
 b_\varepsilon(x)=A_0+\varepsilon B_0\chi(x),\qquad
 Q_\varepsilon(u)=\sum_{x,j,r}b_\varepsilon(x)|\Delta_j u_r(x)|^2.
\end{equation}
Thus $\varepsilon=-1$ has strong weight $H$ on odd source vertices; the
other channel exchanges the two parities. Write $A_\varepsilon$ for the
self-adjoint operator of this form \emph{restricted to $U$}. We do not
assert that the ambient weighted-gradient operator preserves $U$.
The two-channel form and domain are the original direct sum.

For each $i\in\Z_n$ choose $\tau_i\in\{0,\ldots,n-1\}$. The original
nonwrapping cells are
\[
 \{i\}\times I\times J,\quad
 I\in\{[0,\tau_i-1],\{\tau_i\},[\tau_i+1,n-1]\},\quad
 J\in\{\{0\},[1,n-1]\},
\]
with empty intervals omitted. Every spatial component and channel is
observed separately, with
\begin{equation}\label{eq:C}
 (Cu)_{r,B}=\frac{1}{\sqrt{n^3|B|}}\sum_{x\in B}u_r(x).
\end{equation}
No channel rotation, adapted mean or projected row replaces these observations.
Unless stated otherwise, proofs are per channel; direct sums give the same
constants for both channels.

Our quantitative specialization is
\begin{equation}\label{eq:special}
             H=cn^2,\qquad a=c/48,\qquad c>0.
\end{equation}
When $f(i)=i+\tau_i\pmod n$ is a permutation, denote by $T$ the genuine
singleton submap
\begin{equation}\label{eq:T}
             (Tu)_i=n^{-3/2}u(i,\tau_i,0).
\end{equation}
The quadratic cut $\tau_i=2i^2\pmod n$ has this property for dyadic $n$:
$f(i)=f(j)$ implies $(i-j)(1+2(i+j))=0\pmod n$, and the second factor is
odd and hence invertible.

For the endpoint cut $\tau_i=0$, let $C_{\rm end}$ denote all four
original cells on each slice, in every spatial component. Put
\[
 \lambda_{\rm end}(n)=\min_{0\ne w\in U\cap\ker C_{\rm end}}
                              \frac{Q(w)}{\norm w^2},
 \qquad N=n/2,\qquad h_N=\sum_{j=1}^N\frac1j.
\]
For $n\ge16$ this kernel is nonzero: $\dim U=2(n^3-1)>12n$, the
number of scalar observation rows. Define
\begin{align}\label{eq:mainconstants}
 K_n&=\frac{48}{13}h_N+\frac1{6n^2-1},\\
 b_{\rm end}(n)&=\frac{h_N}{16\pi^2}-\frac{17}2
       -32(2\pi+8)^2-\frac{256K_n^2}{48n^2-1}.\notag
\end{align}

\begin{theorem}[Sharp endpoint-kernel order]\label{thm:endpoint}
Under~\eqref{eq:special}, for even $n\ge16$ with $b_{\rm end}(n)>0$,
\begin{equation}\label{eq:endpointsandwich}
 \frac c{K_n}\le\lambda_{\rm end}(n)
 \le\frac{c\pi^2}{2n^2}+\frac{16c}{13b_{\rm end}(n)}.
\end{equation}
In particular $\lambda_{\rm end}(n)=\Theta(c/\log n)$ along even sizes,
with positive comparison constants independent of $n$ and $c$.
The same bounds hold for both original channels and their direct sum.
For all sufficiently large even sizes there is a nonzero real field
annihilating every endpoint observation with $Q(w)<c\norm w^2$.
\end{theorem}

The lower bound is proved in Corollary~\ref{cor:kernel}. The upper
construction, including its realness and all original cell cancellations,
is proved in Section~\ref{sec:endpoint}. No exact asymptotic constant or
small-size threshold is claimed. The theorem concerns this fixed endpoint
cut, not the best canonical partition.
\section{Complete spectrum on the actual coexact domain}\label{sec:spectrum}
Use the unitary convention
\[
 u(x)=n^{-3/2}\sum_k\widehat u(k)\e^{2\pi i k\cdot x/n},\qquad
 g_j(k)=\e^{2\pi i k_j/n}-1.
\]
Complexification is used for calculation; real multiplicities are recovered
without doubling the complex dimension.

\begin{lemma}[The actual image]\label{lem:image}
At $k\ne0$, the Fourier image of $d_1^*$ is $g(k)^\perp$ of dimension two;
at $k=0$ it is zero. Consequently $\dim_\R U=2(n^3-1)$.
\end{lemma}
\begin{proof}
For a skew array $W_{j\ell}=-W_{\ell j}$,
$(d_1^*W)_\ell=\sum_j\overline{g_j}W_{j\ell}$, which is orthogonal to
$g$. Conversely, for $v\perp g$ and $D=|g|^2>0$, set
\[
              W_{j\ell}=\frac{g_jv_\ell-g_\ell v_j}{D}.
\]
Its adjoint image is $v_\ell-g_\ell(g^*v)/D=v_\ell$. At zero all symbols
vanish. The formula respects conjugation at negative frequencies, so real
fields have real preimages. Counting the nonzero frequencies proves the
dimension and excludes harmonic constants.
\end{proof}

Let $q=(n/2,n/2,n/2)$ and consider an unordered pair $\{k,k+q\}$. Put
$h=g(k+q)$, $D=|g|^2$, $D'=|h|^2$ and $\Gamma=g^*h$. Direct calculation gives
\begin{equation}\label{eq:symbols}
 D'=12-D,\qquad \Gamma=2i\sum_j\sin(2\pi k_j/n).
\end{equation}

\begin{theorem}[All restricted eigenvalues]\label{thm:spectrum}
For a nonexceptional pair, both frequencies being nonzero, the four
eigenvalues per channel are
\begin{align}\label{eq:roots}
 \lambda_{1,\pm}&=6A_0\pm\sqrt{A_0^2(D-6)^2+B_0^2|\Gamma|^2},\\
 \lambda_{2,\pm}&=6A_0\pm\sqrt{A_0^2(D-6)^2+
                           B_0^2|\Gamma|^4/(DD')}.\notag
\end{align}
Repeated roots retain their multiplicities. The pair $\{0,q\}$ has only
the $q$ plane and its two eigenvalues are $12A_0$. These blocks exhaust
$U$: their total multiplicity is $4(n^3/2-1)+2=2(n^3-1)$ per channel.
Both original channels have the same spectrum.
\end{theorem}
\begin{proof}
For coefficients $x=\widehat u(k)$, $y=\widehat u(k+q)$, the original
source-weighted form is exactly
\[
 A_0D|x|^2+A_0D'|y|^2+2\varepsilon B_0\Rea(\Gamma x^*y).
\]
Indeed multiplication by $\chi$ shifts frequency by $q$, and summing the
three forward differences gives $g^*h$ with no extra volume factor.
Choose isometries $E,F:\C^2\to\C^3$ onto $g^\perp,h^\perp$. The compressed
block has diagonal $A_0DI_2,A_0D'I_2$ and off-diagonal
$\varepsilon B_0\Gamma E^*F$.

The singular values of $E^*F$ are $1$ and
$\rho=|\Gamma|/\sqrt{DD'}$. The planes have a common unit vector, and their
projections $P_g=I-gg^*/D$, $P_h=I-hh^*/D'$ satisfy
$\tr(P_gP_h)=1+|\Gamma|^2/(DD')$. This proves the two singular values,
including coincident or orthogonally normal planes. Singular-value
decomposition and absorption of the coupling phase give two scalar
two-by-two blocks with coupling moduli $B_0|\Gamma|$ and $B_0|\Gamma|\rho$.
Their eigenvalues are~\eqref{eq:roots}.

At zero the first coefficient is absent by Lemma~\ref{lem:image}; at $q$
the remaining plane has $D'=12$. Changing the channel sign is conjugation
by a sign on the second coefficient. Complex conjugation takes each block
to its negative-frequency block and commutes with the real form, so the
real eigenspace dimensions equal those of the complexification. The count
above is therefore also the real count, not twice that count.
\end{proof}

\begin{proposition}[The whole low space]\label{prop:low}
Under~\eqref{eq:special}, let $L=1_{(0,c)}(A_\varepsilon)U$ and
$F=U\ominus L$. Then
\[
 \dim L=n-2,\qquad A_L\le(c/8)I,\qquad A_F\ge(16c/3)I.
\]
The space $L$ consists exactly of the low branches at nonexceptional
diagonal pairs $k=(t,t,t)$; both transverse polarizations are retained.
\end{proposition}
\begin{proof}
First keep only the strong source weight, setting $a=0$. Put
$S=\sum_j\e^{2\pi i k_j/n}$. Formula~\eqref{eq:roots} gives the attained
smallest coexact pair eigenvalue $H(3-|S|)$. It vanishes exactly when the
three phases agree. At such a nonexceptional diagonal pair both couplings
have the same singular value and the zero multiplicity is two. The
exceptional pair contributes no zero mode. Thus the actual strong kernel
has dimension $n-2$ per channel.

If the residues are not all equal, at least two pairwise differences are
nonzero. The identity
\[
 9-|S|^2=4\sum_{j<\ell}\sin^2\!\left(\frac{\pi(k_j-k_\ell)}n\right)
\]
and division by $3+|S|\le6$ show
$H(3-|S|)\ge(4H/3)\sin^2(\pi/n)\ge16H/(3n^2)$.
Here $\sin(\pi/n)\ge2/n$ is the concavity chord on $[0,\pi/2]$.
Adding the positive weak form preserves this lower bound.

At a diagonal nonexceptional pair, $\theta=2\pi t/n$, the two roots,
each with multiplicity two, are
\[
 \mu_\pm=3(H+a)\pm3\sqrt{(H+a)^2-4Ha\sin^2\theta}.
\]
They satisfy $0<\mu_-\le6a=c/8$ and $\mu_+\ge6H$. The exceptional
eigenvalue is $6(H+a)\ge6H$. There are $n/2-1$ nonexceptional diagonal
pairs, proving the count and the complete separation. In particular no
other low direction has been silently omitted.
\end{proof}

\section{Original rows and a paid range theorem}\label{sec:range}
For a cell $B=\{i\}\times I\times J$, direct substitution gives
\begin{equation}\label{eq:Fourierrow}
 (Cu)_{r,B}=\frac1{n^3\sqrt{|I||J|}}\sum_k\widehat u_r(k)
       \e^{2\pi i k_1i/n}
       \left(\sum_{j\in I}\e^{2\pi i k_2j/n}\right)
       \left(\sum_{\ell\in J}\e^{2\pi i k_3\ell/n}\right).
\end{equation}
Each pair member has its own three factors. The zero coefficient is absent
on $U$. Finite sums, rather than a geometric quotient at zero, fix the
normalization without removable singularities.

Let $N_d$ be the actual diagonal Fourier space: its fields have the form
$u(x)=v(x_1+x_2+x_3)$ with mean-zero $v$ and $v_1+v_2+v_3=0$. This follows
from Lemma~\ref{lem:image}. Thus $\dim N_d=2(n-1)$ and $L\subset N_d$.
Put $O=U\ominus N_d$. The following result permits every original low
observation to be cancelled by actual high modes, including row dependencies.

\begin{theorem}[Range equality and explicit lift]\label{thm:range}
For every original cut, whether or not $f$ is a permutation,
\[
       \ran C_O=\ran(C|_U)=\ran C_F.
\]
There is a linear $R_d:N_d\to O$ such that
\[
 C_OR_d=-C_{N_d},\qquad
 \norm{R_d}^2\le K_n^{\rm lift}:=\frac{n^5}{16(n-3)}.
\]
\end{theorem}
\begin{proof}
For any original row coefficient $\alpha$, put $g=C_{\rm amb}^*\alpha$,
where the ambient map uses precisely~\eqref{eq:C}. For fixed $x_1,x_2$,
$g$ is constant in $x_3$ on $1,\ldots,n-1$, because that is a single long
third-coordinate cell. Let $v=P_Ug$. Since $U^\perp=\ker d_1$,
$d_1g=d_1v$. Every component of this curl is independent of $x_3$ on
$1,\ldots,n-2$; its third difference therefore vanishes on $1,\ldots,n-3$.

If $C_O^*\alpha=0$, then $v\in N_d$. Each component of $d_1v$ depends
only on $s=x_1+x_2+x_3$. Equality of its values at $x_3=1,2$, for all
$x_1,x_2$, forces equal values at adjacent residues of $s$. It is constant,
and its sum is zero because it is a periodic curl. Hence $d_1v=0$ and
$v\in U\cap U^\perp=0$. We have proved equality of the annihilators of
$\ran C_O$ and $\ran(C|_U)$. Finite-dimensional orthogonal complementation
gives the first range equality; $O\subset F\subset U$ gives the other.

For the quantitative claim decompose $v=d+o\in N_d\oplus O$. On the slab
$1\le x_3\le n-3$, $\Delta_3d_1d=-\Delta_3d_1o$. The first expression
depends only on $s$, so its squared norm is equally distributed among
the $n$ third-coordinate slices. Since the curl symbol has squared norm
at most $12$ and $|g_3|^2\le4$,
\[
 \norm{\Delta_3d_1d}^2\le\frac{48n}{n-3}\norm{o}^2.
\]
At a nonzero diagonal frequency $z=\e^{2\pi it/n}$ the same squared symbol
on $d$ is $3|z-1|^4$. Using $|z-1|^2\ge16/n^2$ gives
$\norm{\Delta_3d_1d}^2\ge768n^{-4}\norm d^2$. Consequently
\begin{equation}\label{eq:anglelift}
       \norm{C_{N_d}^*\alpha}^2
       \le K_n^{\rm lift}\norm{C_O^*\alpha}^2.
\end{equation}
Define $S(C_O^*\alpha)=C_{N_d}^*\alpha$. Inequality~\eqref{eq:anglelift}
makes this well-defined on its range, with norm at most $\sqrt{K_n^{\rm lift}}$.
Extend it by zero on the orthogonal complement. Then $C_OS^*=C_{N_d}$,
and $R_d=-S^*$ proves the claim. This is a proved factorization, not an
assumed inverse or a full-row-rank assertion.
\end{proof}

\section{The exact shifted cancellation criterion}\label{sec:cost}
Under~\eqref{eq:special}, set $B=A_F-cI$, $D_L=cI-A_L$ and
\[
 M=C_FB^{-1}C_F^*,\qquad S_C=\ran(C|_U).
\]
Proposition~\ref{prop:low} gives $B\ge13c/3$ and $7c/8\le D_L<cI$.
Theorem~\ref{thm:range} implies $\ker M=S_C^\perp$. Thus $M$ is invertible
on $S_C$, and $M^\dagger$ denotes that inverse extended by zero.

\begin{theorem}[Exact cost, with dependent rows]\label{thm:cost}
For every $l\in L$,
\begin{align}\label{eq:mincost}
 \min_{f\in F:\ C_Ff=-C_Ll}\ip f{Bf}
       &=\ip{C_Ll}{M^\dagger C_Ll},\\
 f_*&=-B^{-1}C_F^*M^\dagger C_Ll.\notag
\end{align}
Consequently $Q(w)\ge c\norm w^2$ for all $w\in U\cap\ker C$ if and only if
\begin{equation}\label{eq:criterion}
                  C_L^*M^\dagger C_L-D_L\ge0\quad\hbox{on }L.
\end{equation}
Its dimension is $n-2$ per original channel, or $2(n-2)$ on their direct sum.
\end{theorem}
\begin{proof}
Let $y=C_Ll\in S_C$. Then $C_Ff_*=-MM^\dagger y=-y$. For every other
feasible $f=f_*+h$, $C_Fh=0$ and
$\ip h{Bf_*}=-\ip{C_Fh}{M^\dagger y}=0$. Expanding the square proves
the minimum. For $w=l+f$ the spectral splitting gives
\[
        Q(w)-c\norm w^2=-\ip l{D_Ll}+\ip f{Bf}.
\]
Use the attained minimum in both directions. Every $l$ is feasible by
Theorem~\ref{thm:range}, so no extra range assumption hides a bad direction.
\end{proof}

The lift also pays an explicit but inadequate upper bound:
$\ip{R_dl}{BR_dl}\le(12H-c)K_n^{\rm lift}\norm l^2$, because
$A\le12HI$ and $R_dl\in O\subset F$. This grows like $cn^6$ and does not
decide~\eqref{eq:criterion}.
For later use the corresponding dual identity is
\begin{equation}\label{eq:dual}
 \ip y{M^\dagger y}
   =\sup_\alpha\{2\Rea\ip\alpha y-\ip\alpha{M\alpha}\},\qquad y\in S_C.
\end{equation}
It follows by completing the square on $S_C$; components in $S_C^\perp$
pair and respond to zero.

\begin{lemma}[Singleton frame]\label{lem:frame}
If $f$ is a permutation, $T^*T=n^{-5}I$ as a form on all of $N_d$.
In particular this holds on $L$ and on the strong kernel.
\end{lemma}
\begin{proof}
For $u(x)=v(s)$, $\norm u^2=n^2\sum_s|v(s)|^2$, whereas
$\norm{Tu}^2=n^{-3}\sum_i|v(f(i))|^2=n^{-3}\sum_s|v(s)|^2$.
\end{proof}

\section{Whole singleton response and an invisible-field bound}\label{sec:whole}
Define $R=P_FB^{-1}P_F$ extended by zero and $M_T=TRT^*$. Put
$N=n/2$, $h_N=\sum_{j=1}^N1/j$ and
\begin{equation}\label{eq:Kn}
                   K_n=\frac{48}{13}h_N+\frac1{6n^2-1}.
\end{equation}

\begin{theorem}[Finite whole-response bounds]\label{thm:whole}
For every permutation cut and either original channel,
\begin{equation}\label{eq:whole}
 \frac{h_N}{24\pi^2cn^5}\le\norm{M_T}\le\frac{K_n}{cn^5}.
\end{equation}
No fixed constant bounds $cn^5\norm{M_T}$ for all such sizes. The same
blanket scalar bound fails for the entire original-row response $M$.
\end{theorem}
\begin{proof}
Write $k=(t+r,t+s,t)$ and
$\delta(r,s)=3-|1+\e^{2\pi ir/n}+\e^{2\pi is/n}|$.
For $g=T^*\alpha$,
\[
 \widehat g(t+r,t+s,t)=n^{-3}\sum_i\alpha_i
   \e^{-2\pi i\{tf(i)+ri+s\tau_i\}/n}.
\]
Since $f$ is a permutation, orthogonality in $t$ gives the exact identity
\begin{equation}\label{eq:fibre}
       \sum_t|\widehat g(t+r,t+s,t)|^2=n^{-5}\norm\alpha^2.
\end{equation}
The pair shift $q$ changes only $t$ by $n/2$. Every transverse $(r,s)$
therefore contains exactly $n/2$ unordered pairs, and only $(0,0)$ is
diagonal. Off that fibre, Proposition~\ref{prop:low} and its strong-form
calculation give
\[
 (A-cI)^{-1}\le\frac1{cn^2\delta-c}I
             \le\frac{16}{13cn^2\delta}I.
\]
The last step uses $cn^2\delta\ge16c/3$. On the diagonal high space the
bound is $1/[c(6n^2-1)]$. Projections cannot increase the input norm.
Thus~\eqref{eq:fibre} gives
\begin{equation}\label{eq:responseupper}
 \ip\alpha{M_T\alpha}\le\frac{\norm\alpha^2}{cn^5}
 \left\{\frac{16}{13n^2}\sum_{(r,s)\ne0}\frac1\delta
                           +\frac1{6n^2-1}\right\}.
\end{equation}
For unique centered residues in $[-N,N-1]^2$, the sine identity in the
proof of Proposition~\ref{prop:low} implies
\[
 \delta\ge\frac23\{\sin^2(\pi r/n)+\sin^2(\pi s/n)\}
          \ge\frac8{3n^2}(r^2+s^2).
\]
At most $8j$ centered sites have $\max(|r|,|s|)=j$. Therefore
\begin{equation}\label{eq:latticesums}
 \sum_{(r,s)\ne0}\frac1\delta\le3n^2h_N,
 \qquad
 \sum_{(r,s)\ne0}\frac1\delta\ge\frac{n^2h_N}{4\pi^2}.
\end{equation}
For the lower inequality retain $r=j$, $1\le s\le j$, $1\le j\le N$.
These are distinct residues and
$\delta\le2-\cos(2\pi r/n)-\cos(2\pi s/n)\le4\pi^2j^2/n^2$.
The upper inequality and~\eqref{eq:responseupper} prove the upper norm.

For the lower norm, even-sum translations preserve $U$, the source form,
$F$ and $R$. They are transitive on each vertex parity. The spatial trace
of $R(x,x)$ therefore has two values. The sampled vertex parity is that
of $f(i)$, so exactly half the singleton points have each parity. Compared
with $n^3/2$ vertices of each parity, the row normalization gives
\begin{equation}\label{eq:tracewhole}
                  \tr M_T=n^{-5}\tr_F B^{-1}.
\end{equation}
On each non-diagonal pair one genuine strong-minimizing polarization has
energy $cn^2\delta$. The weak form is at most $12a=c/4$ on a unit vector.
Rayleigh's principle supplies a full eigenvalue
$\lambda\le cn^2\delta+c/4$. It belongs to $F$; hence
$(\lambda-c)^{-1}\ge(cn^2\delta)^{-1}$. There are exactly $n/2$ such
pairs per transverse fibre. Consequently
\[
 \tr_F B^{-1}\ge\frac1{2cn}\sum_{(r,s)\ne0}\frac1\delta
              \ge\frac{nh_N}{8\pi^2c}.
\]
The largest eigenvalue of a PSD operator on $3n$ rows is at least its
trace divided by $3n$. This proves the lower half of~\eqref{eq:whole}.
Since $h_N\ge\log(N+1)$, the scaled norm diverges. Finally, inserting
singleton coefficients by zeros is an isometry $J$ with $T=J^*C$, so
$M_T=J^*MJ$ and $\norm M\ge\norm{M_T}$. No independent-row assumption
is used in this compression argument.
\end{proof}

\begin{corollary}[A paid logarithmic kernel bound]\label{cor:kernel}
For every $w\in U\cap\ker C$ under a permutation cut,
\begin{equation}\label{eq:kernel}
 Q(w)\ge\frac c{K_n}\norm w^2
             \ge\frac c{4(1+\log n)}\norm w^2.
\end{equation}
The same coefficient holds for the original two-channel direct sum.
\end{corollary}
\begin{proof}
For $l\in L$, insert $\alpha=Tl$ into the singleton coordinates of the full
rowspace. Its pairing with $y=C_Ll$ is $n^{-5}\norm l^2$, and
Theorem~\ref{thm:whole} bounds its response by
$K_n\norm l^2/(cn^{10})$. Optimize~\eqref{eq:dual} along $t\alpha$ at
$t=cn^5/K_n$ to obtain
$\ip{C_Ll}{M^\dagger C_Ll}\ge(c/K_n)\norm l^2$.
For an invisible $w=l+f$, the high part is feasible and
\[
 Q(w)=\ip l{A_Ll}+\ip f{Bf}+c\norm f^2
       \ge(c/K_n)\norm l^2+c\norm f^2
       \ge(c/K_n)\norm w^2,
\]
since $K_n>1$. Put $L_0=1+\log n$. The inequalities
$h_N\le1+\log(n/2)\le L_0$ and
$1/(6n^2-1)\le1/95<(4/13)L_0$ give $K_n<4L_0$.
The proofs are channelwise, so summation preserves the coefficient.
\end{proof}

\begin{remark}[Coefficient and domain comparison]\label{rem:comparison}
The coefficient $c/K_n$ strictly exceeds the prior working-paper benchmark
$c/[5(1+\log n)]$~\cite{Eriksson274} under the same permutation-cut hypothesis. That estimate allowed
ambient fields;~\eqref{eq:kernel} is on
the actual coexact domain with a permutation cut. The numerical improvement
does not extend the larger ambient domain, change logarithmic order, or
establish the target $c$. In particular $c/K_n$ is inadequate to dominate
$D_L\ge7c/8$ in~\eqref{eq:criterion}.
\end{remark}

\section{Logarithmic response persists on useful low data}\label{sec:useful}
Let $E=T_LL$ and let $P_E$ be the ordinary singleton row-space projection.
Lemma~\ref{lem:frame} gives $P_E=n^5T_LT_L^*$ and $\dim E=n-2$.

\begin{theorem}[Useful-space obstruction]\label{thm:useful}
For every even $n\ge4$ and permutation cut,
\begin{equation}\label{eq:useful}
 \frac{h_N/(16\pi^2)-1/2-4K_n/(48n^2-1)}{cn^5}
 \le\norm{P_EM_TP_E}\le\frac{K_n}{cn^5}.
\end{equation}
For even $n\ge16$ the lower bound implies
\begin{equation}\label{eq:simpleuseful}
        cn^5\norm{P_EM_TP_E}\ge\frac{h_N}{32\pi^2}-1.
\end{equation}
Thus a uniform response bound fails even for unaugmented tests $\alpha=T_Ll$.
The lower expressions need not be positive at small sizes.
\end{theorem}

\subsection{A finite component Green estimate}
Let $K$ be the actual strong kernel. For strong weight on odd vertices,
the strong scalar graph has $N$ components, the paired layers
$(s=2j-1,s=2j)$, $s=x_1+x_2+x_3$. On each layer use coordinates
$y=(x_1,x_2)$; $x_3$ is determined. The three strong edges from an odd
vertex have target coordinates $y,y+e_1,y+e_2$ in the even layer. The
unit-weight scalar component Laplacian $D$ is degree three on $2n^2$
vertices, with transverse Fourier matrix
\[
 \begin{pmatrix}3&-S(r,s)\\-\overline{S(r,s)}&3\end{pmatrix},\qquad
 S(r,s)=1+\e^{2\pi ir/n}+\e^{2\pi is/n}.
\]
Its eigenvalues are $\delta=3-|S|$ and $\delta_+=3+|S|$. Only $\delta(0,0)$
vanishes, proving connectedness. Interchanging source parity gives the
same graph with paired layers $(2j,2j+1)$.

\begin{lemma}[Component Green bounds]\label{lem:green}
For $G=D^\dagger$, the diagonal value $g_0$ is constant and
\[
                   g_0\ge\frac{h_N}{8\pi^2},\qquad
                   G(x,y)\ge-1\quad\hbox{for all }x,y.
\]
\end{lemma}
\begin{proof}
Translations within layers and layer exchange with reflection act
transitively on vertices. The diagonal spectral formula is
\[
 g_0=\frac1{2n^2}\left\{
       \sum_{(r,s)\ne0}\frac1\delta+\sum_{r,s}\frac1{\delta_+}\right\}.
\]
The lower lattice sum in~\eqref{eq:latticesums} gives the diagonal bound.
For the off-diagonal bound use the finite heat kernel $p_t=\exp(-tD)$.
It is entrywise nonnegative, since $D=3I-\mathrm{Adj}$ and
$\exp(t\mathrm{Adj})$ has a nonnegative power series. Hence the integral
representation
\[
 G(x,y)=\int_0^\infty\{p_t(x,y)-(2n^2)^{-1}\}\,dt
\]
is bounded below by $-1/2$ on $0\le t\le n^2$. For the tail,
$p_t-P_{\rm const}$ is PSD with constant diagonal. Cauchy--Schwarz bounds
each entry by that diagonal, its trace divided by $2n^2$. The absolute
tail is therefore at most
\[
            \frac1{2n^2}\sum_{\lambda\ne0}
                         \frac{\e^{-n^2\lambda}}\lambda.
\]
For the $\delta$ branch use $\delta\ge8(r^2+s^2)/(3n^2)$ and the at-most
$8j$ shell count. Its contribution is at most
\[
 \frac32\sum_{j\ge1}\frac{\e^{-(8/3)j^2}}j
       \le\frac32\sum_{j\ge1}(1/8)^j=\frac3{14}.
\]
Here $\e^{8/3}>8$ follows already from four nonnegative Taylor terms.
The $n^2$ high-branch eigenvalues are at least $3$, giving a contribution
at most $\e^{-3n^2}/6\le1/6$. Thus
$G(x,y)\ge-1/2-3/14-1/6=-37/42\ge-1$. No positivity of the Green function
or an isolated low band was assumed.
\end{proof}

\subsection{The scalar component inside the actual coexact space}
Put $V=(1,1,1)^\perp$. The strong kernel $K$ consists exactly of component
constants $v_j\in V$ with $\sum_jv_j=0$. Such fields have mean-zero
diagonal Fourier coefficients in $V$, so Lemma~\ref{lem:image} puts them
in $U$. Their dimension $2(N-1)=n-2$ equals the full strong-kernel count.
Every permutation singleton map selects one point in each $s$ layer,
hence one point $o_j,e_j$ in each layer of each component. The space
$E_0=T_KK$ consists of equal pair values in $V$, with component mean zero.

The next construction prevents an ambient scalar inverse from being
mistaken for an actual coexact inverse. At transverse $(r,s)\ne0$, let
$z=(\e^{2\pi ir/n},\e^{2\pi is/n},1)$ and
$W_{r,s}=\{(1,1,1),z\}^\perp$, a line in $V$. For frequencies
$k=(t+r,t+s,t)$ and $k+q$, both $g(k)$ and $g(k+q)$ are linear combinations
of $z$ and $(1,1,1)$. Thus $W_{r,s}$ lies in both actual coexact planes,
and is independent of $t$. Its two coefficient directions form an invariant
first-polarization subspace of the strong block, with eigenvalues
$H\delta,H\delta_+$. Retaining this line in the transverse Fourier
transform of each two-layer component gives actual coexact fields. The
projector onto the line depends only on transverse frequency, so it acts
inside each component and does not mix components.

Let $R_0$ be this selected strong inverse, zero on the entire diagonal
transverse fibre. It is $H^{-1}$ times the component scalar inverse with
spatial projector $P_{W_{r,s}}$. On every non-diagonal actual block,
\[
           0<B=A-cI\le A_H-3cI/4\le A_H,
\]
because $0\le A_a\le cI/4$. Positive matrix inversion gives
$B^{-1}\ge A_H^{-1}\ge R_0$. The last comparison keeps an invariant
polarization of the strong inverse. All such blocks lie in $F$, while
the remaining diagonal part of $R$ is positive. Extending by zero gives
$R\ge R_0$ globally. This comparison never inverts the strong kernel and
never compares inverses on the rotated diagonal low modes.

For the normalized pair vector
$p_j=(\delta_{o_j}+\delta_{e_j})/\sqrt2$, taking spatial trace against
$P_V$ removes the polarization orientation:
$\tr(P_VP_{W_{r,s}})=1$. The omitted diagonal transverse high branch
would contribute zero: its layer vector is $(1,-1)/\sqrt2$, constant in
transverse coordinates, and pairs to zero with $p_j$. The scalar zero
branch is already zero in $G$. Therefore the local useful-pair trace is
\begin{equation}\label{eq:localtrace}
 n^{-3}H^{-1}\ip{p_j}{Gp_j}
   =\frac{g_0+\Rea G(o_j,e_j)}{cn^5}
   \ge\frac{h_N/(8\pi^2)-1}{cn^5}.
\end{equation}
There is no factor two from $V$: only one common polarization was retained
per transverse frequency. The other positive contributions are omitted in
a lower estimate, not removed from the original model.

The component mean-zero projection is $I_N-N^{-1}11^*$, with diagonal
$1-1/N$. Since $R_0$ is componentwise, summing~\eqref{eq:localtrace}
and using $R\ge R_0$ proves
\begin{equation}\label{eq:strongtrace}
 \tr(P_{E_0}M_T)\ge\frac{N-1}{cn^5}
                            \left(\frac{h_N}{8\pi^2}-1\right).
\end{equation}
This is a useful strong-kernel compression trace, not the blanket trace
in~\eqref{eq:tracewhole}. Arbitrary separation of $o_j,e_j$ is allowed.

\subsection{Positive-weight rotation and conclusion}
We must not identify $L$ with $K$. On each nonexceptional diagonal pair
the strong scalar block has eigenvalues $0,6H$, tensored with $V$.
The weak perturbation has norm at most $c/4$, and the true low eigenvalue
is at most $c/8$. For its unit eigenvector $l=l_0+l_1$ in the strong
low/high basis, the high equation gives
\[
 (6H-\lambda)l_1=-P_{\rm high}A_al,\qquad
 \norm{l_1}\le\frac{c/4}{6cn^2-c/8}
                     =:\epsilon_n=\frac2{48n^2-1}.
\]
The norm difference of two rank-one projections is the sine of their
angle, as follows by writing the projections in the strong orthonormal
basis. Tensoring with $V$ and taking the orthogonal sum over pairs gives
$\norm{P_L-P_K}\le\epsilon_n$. The exceptional and non-diagonal low
projections are both zero.

By Lemma~\ref{lem:frame}, $n^{5/2}T$ is an isometry on all $N_d$, so
$\norm{P_E-P_{E_0}}\le\epsilon_n$. Both ranks are $d=n-2$; their
difference has rank at most $2d$ and trace norm at most $2d\epsilon_n$.
Using Theorem~\ref{thm:whole},
\[
 |\tr((P_E-P_{E_0})M_T)|\le\frac{2d\epsilon_nK_n}{cn^5}.
\]
Combine this with~\eqref{eq:strongtrace}, divide by $d=2(N-1)$ and use
PSD norm at least average eigenvalue. The result is the lower half of
\eqref{eq:useful}; the upper half follows by orthogonal compression.

For $n\ge16$, $D_n=48n^2-1\ge12287$ and
$(192/13)/D_n\le1/(32\pi^2)$, since $\pi<4$ implies
$6144\pi^2<98304<13\cdot12287$. Also
$4/[(6n^2-1)D_n]<1/2$. Formula~\eqref{eq:Kn} therefore gives
$4K_n/D_n\le h_N/(32\pi^2)+1/2$, which proves~\eqref{eq:simpleuseful}.
Since $h_N\ge\log(N+1)$, the scaled useful norm diverges.

The derivation is valid after complexification; the original response and
useful projection are real, with the same traces and eigenvalues. A
maximizing direction can be chosen real. Moreover the isometry
$n^{5/2}T_L:L\to E$ proves the exact identity
\[
 \sup_{l\ne0}\frac{cn^{10}\ip{T_Ll}{M_TT_Ll}}{\norm l^2}
                   =cn^5\norm{P_EM_TP_E}.
\]
This is a large dual response, not a physical Rayleigh quotient.

\section{Sharp loss for all endpoint observations}\label{sec:endpoint}
We now fix $\tau_i=0$ for every $i$ and write $C_{\rm end}$ for the
\emph{full} map~\eqref{eq:C}. Its cells on each first-coordinate slice are
the four products of $\{0\}$ and $[1,n-1]$. The singleton map is evaluation
on $(i,0,0)$. Throughout this section $n\ge16$ is even. We prove
Theorem~\ref{thm:endpoint} by retaining a low-frequency sector while
cancelling every endpoint cell, including the long rectangles.

\subsection{A selected response and its cross terms}
Let $S=n^{3/2}T$ be unnormalized evaluation and $W=V^n$ in singleton
row space. Using the actual common-polarization inverse $R_0$ constructed
in Section~\ref{sec:useful}, define
\[
                 \mathcal K=S(HR_0)S^*.
\]
It annihilates $W^\perp$. On $W$ it is a direct sum of two-point blocks,
one for each strong component. Write these real blocks as
\[
 \mathcal K_j=\begin{pmatrix} A_0'&B_j'\\(B_j')^t&A_0'\end{pmatrix}.
\]
The primes distinguish these matrices from the scalar source weights.
For $\theta=(r,s)\ne0$, put $z=(\e^{ia},\e^{ib},1)$,
$a=2\pi r/n$, $b=2\pi s/n$, and let $P(\theta)$ project onto
$\{1,z\}^\perp$. With $S_\theta=1+\e^{ia}+\e^{ib}$ and
$\Lambda_\theta=9-|S_\theta|^2$, the component Fourier formula is
\[
 A_0'=n^{-2}\sum_{\theta\ne0}\frac{3P(\theta)}{\Lambda_\theta},
 \qquad
 B_j'=n^{-2}\sum_{\theta\ne0}
        \frac{\zeta_j(\theta)S_\theta P(\theta)}{\Lambda_\theta},
 \qquad |\zeta_j|=1.
\]
Conjugating $S_\theta$ when exchanging layers gives the same bounds.
The evaluation phase $\zeta_j$ records the original two sample points.

\begin{lemma}[Polarization average and bounded cross block]\label{lem:cross}
On $W$, $\mathcal K\le K_n I$ and
\[
 \mathcal K\ge\frac{1-n^{-2}}{12}I.
\]
After changing to pair-plus and pair-minus coordinates, its local blocks
$X_j,Y_j,Z_j$ obey
\[
 \tr X_j\ge \frac{h_N}{8\pi^2}-1,
 \qquad \norm{Z_j}\le2\pi.
\]
\end{lemma}
\begin{proof}
The exact polarization average is
\begin{equation}\label{eq:polaverage}
 n^{-2}\sum_{\theta\ne0}P(\theta)=\frac{1-n^{-2}}2P_V.
\end{equation}
Indeed conjugate frequencies make the sum real. Permute the three
coordinates of $z$ and divide by its last coordinate: this bijects the
nonzero transverse frequencies and conjugates $P$ by that spatial
permutation. A real matrix commuting with all three-coordinate
permutations has constant diagonal and constant off-diagonal entries.
It annihilates $1$, so is scalar on $V$. Its trace is $n^2-1$, proving
\eqref{eq:polaverage}. Since the two eigenvalues of each nonzero component
matrix lie in $(0,6]$, its inverse is at least $I/6$.
Evaluation on one point in each layer and~\eqref{eq:polaverage} give the
floor. The upper bound follows from $R_0\le R$ and
Theorem~\ref{thm:whole}, since $Hn^3=cn^5$.

The plus block is $X_j=A_0'+(B_j'+(B_j')^t)/2$. Its spatial trace is
$g_0+G(o_j,e_j)$, by~\eqref{eq:localtrace} after removing its normalization.
The zero transverse optical branch contributes zero to the plus vector.
Lemma~\ref{lem:green} proves the trace bound.

The cross block is $Z_j=((B_j')^t-B_j')/2$ up to a sign. To bound it,
choose centered angles $a,b\in[-\pi,\pi)$ and set $\rho=|a|+|b|$.
A spanning vector for the polarization line is
\[
 w=(\e^{-ib}-1,\ 1-\e^{-ia},\ \e^{-ia}-\e^{-ib}),
 \qquad |w|^2=\Lambda_\theta.
\]
Its entries have real amplitudes times a common factor $-2i$ and phases
$\e^{-ib/2},\e^{-ia/2},\e^{-i(a+b)/2}$. Every phase difference is at
most $\rho/2$. The Frobenius norm of the entrywise imaginary part of
$ww^*/|w|^2$ is therefore at most $\rho/2$. Because $B_j'$ is real,
transposition gives
\[
 \norm{Z_j}\le\frac32 n^{-2}
                   \sum_{\theta\ne0}\frac{\rho}{\Lambda_\theta}.
\]
The chord bound gives $\Lambda_\theta\ge4(a^2+b^2)/\pi^2$.
On a centered-residue shell $m=\max(|r|,|s|)$,
$\rho/(a^2+b^2)\le n/(\pi m)$; there are at most $8m$ sites and
$N$ shells. Consequently $n^{-2}\sum\rho/\Lambda_\theta\le\pi$.
Thus $\norm{Z_j}\le3\pi/2\le2\pi$. The phase $\zeta_j$ has modulus
one, so no bound on the distance between sample points was assumed.
\end{proof}

\subsection{Removing line averages inside the original domain}
Let $\Pi_j$ be the ambient Fourier projector onto $k_j\in\{0,N\}$,
for $j=2,3$, and put
\[
 \Pi=\Pi_2+\Pi_3-\Pi_2\Pi_3,\qquad \Xi=I-\Pi,
 \qquad R'=\Xi R\Xi,\qquad R_0'=\Xi R_0\Xi.
\]
These are actual frequency projectors, not projections of observation rows.
The Fourier description of $U$ shows that they preserve $U$. They retain
or remove a whole pair $\{k,k+q\}$, so commute with the restricted $A$.
All low pairs are diagonal with coordinates outside $\{0,N\}$, hence
$\Xi|_L=I$. Therefore $R'$ is precisely the shifted inverse on
$F'=\Xi F$, and $0\le R_0'\le R'\le R$.
Every field in $\ran\Xi$ has zero ordinary and alternating full-line
sums in both the second and third coordinates.

\begin{lemma}[A floor after removal]\label{lem:floorremoved}
Set $\mathcal K'=S(HR_0')S^*$. On $W$,
\begin{equation}\label{eq:removedfloor}
              I/32\le\mathcal K'\le\mathcal K\le K_n I.
\end{equation}
Both selected responses annihilate $W^\perp$, while the actual response
$\widehat M'=cn^5TR'T^*$ satisfies
$\mathcal K'\le\widehat M'\le K_nI$ on the full row space.
\end{lemma}
\begin{proof}
Let $Q_0$ be the support projector of $R_0$: it retains the common
polarization on both layers for each nonzero transverse frequency.
Then $HR_0\ge Q_0/6$. In the original dictionary
$k=(t+r,t+s,t)$, each $\Pi_j$ retains a whole pair for every transverse
frequency, so commutes with $Q_0$ and $R_0$.
Equation~\eqref{eq:polaverage} and one sample in each graph layer give
$SQ_0S^*=(1-n^{-2})I_W/2$. Also $S\Pi_jS^*=2I/n$ on the full
spatial row space: the band kernel vanishes unless the other two
coordinates coincide, and distinct selected points have distinct $x_1$.
Its coincident diagonal is $2/n$. Since $0\le\Pi Q_0\le\Pi_2+\Pi_3$,
\[
 \mathcal K'\ge \frac16S\Xi Q_0\Xi S^*
 \ge\frac{1-8/n-n^{-2}}{12}I_W\ge\frac1{32}I_W.
\]
The scalar expression increases for $n\ge16$ and the last inequality
already holds at $16$. All upper comparisons follow by positivity and
Theorem~\ref{thm:whole}. In particular the floor has been proved after
removal; it does not follow from a trace estimate alone.
\end{proof}

\subsection{The lowest sector loses only a bounded response}
Translation by two in $x_1$ preserves the endpoint sample line and all
operators just defined. The complex row sector $\mathcal H_1$ consists of
\[
 \alpha_i=n^{-1/2}\bigl(a\e^{2\pi i i/n}
                  +b\e^{2\pi i(1+N)i/n}\bigr),\qquad a,b\in\C^3.
\]
Its squared norm is $|a|^2+|b|^2$, and its intersection with $W$ has
dimension four. In paired-component coordinates this is the nonconstant
character $\e^{2\pi i j/N}$. The space $E\cap\mathcal H_1$ has dimension
two: it comes exactly from the diagonal low pair $(1,1,1)$ and its
$q$ partner, with both polarizations. Rationalizing its eigenvalue gives
\begin{equation}\label{eq:lowestenergy}
 \mu_-(1)=\frac{12Ha\sin^2(2\pi/n)}
 {H+a+\sqrt{(H+a)^2-4Ha\sin^2(2\pi/n)}}
 \le6a\sin^2(2\pi/n)\le\frac{c\pi^2}{2n^2}.
\end{equation}
The denominator is at least $2H$, since the square root is at least $H-a$.

\begin{lemma}[Sector loss]\label{lem:sectorloss}
For $\Delta=\mathcal K-\mathcal K'\ge0$,
$\norm{\Delta|_{\mathcal H_1}}\le8$.
\end{lemma}
\begin{proof}
For $g=S^*\alpha$, its unitary torus Fourier coefficients are $a/n$ at
$k_1=1$ and $b/n$ at $k_1=1+N$, independently of $k_2,k_3$.
Represent each unordered checkerboard pair uniquely by $k_1=1$.
Its datum has squared norm $(|a|^2+|b|^2)/n^2$, with no extra pair factor.
On a removed pair at least one of $k_2,k_3$ is $0$ or $N$; this pair
cannot be diagonal. The norm of $HR_0$ there is at most $1/\delta$.
For $k_3$ fixed choose centered residues $r=1-k_3$ and $t=k_2-k_3$.
Then $|r|\ge1$, $t$ runs over all residues, and
$\delta\ge8(r^2+t^2)/(3n^2)$ by the sine bound above. The same argument
applies with $k_2$ fixed, after exchanging the coordinates.
Each of the four removed lines contributes at most
\[
 \frac1{n^2}\sum_t\frac1\delta
 \le\frac38\sum_{t\in\Z}\frac1{1+t^2}
 \le\frac38(1+2\cdot2)=\frac{15}8.
\]
Here $\sum_{t\ge1}t^{-2}\le1+\int_1^\infty x^{-2}\,dx=2$.
Counting intersections twice only enlarges the positive bound.
The four-line total is $15/2\le8$. Polarization projections decrease
the quadratic response, and commutation identifies the removed response
with exactly $\Delta$. Restriction to $W$ preserves the bound.
\end{proof}

\subsection{Optimizing all complementary singleton coefficients}
For positive definite $G$ on $W$ and a subspace $E'\subset W$, write
\[
 \mathcal S_{E'}(G)=(P_{E'}G^{-1}P_{E'}|_{E'})^{-1}.
\]
This is the minimum energy with fixed $P_{E'}\alpha=y$: solving the
multiplier equation $G\alpha\in E'$ gives both this formula and
attainment. It is not the inverse of the direct compression of $G$.

For the endpoint cut the selected two points in each component have
graph coordinates $(2j-1,0),(2j,0)$. Their relative displacement is the
same in every component. Thus $\mathcal K$ has the same local block in
each component, and the bounds of Lemma~\ref{lem:cross} hold in
$\mathcal H_1\cap W$. The strong useful space there is the pair-plus
copy of $V$. Its mean-zero condition is automatic in this nonconstant
character; no additional common-mean minimization is needed.

Write the blocks of $\mathcal K'$ in pair-plus/minus coordinates as
$X',Y',Z'$. Lemmas~\ref{lem:floorremoved} and~\ref{lem:sectorloss} give
\[
 \tr X'\ge\frac{h_N}{8\pi^2}-17,\qquad
 \norm{Z'}\le2\pi+8,\qquad Y'\ge I/32.
\]
Minimizing the minus variable gives the Hermitian Schur complement
$S_0'=X'-Z'(Y')^{-1}(Z')^*$. Since $V$ has dimension two,
\begin{equation}\label{eq:sectortrace}
 \tr S_0'\ge\frac{h_N}{8\pi^2}-17-64(2\pi+8)^2.
\end{equation}
All adjoints in this complex sector are Hermitian.

The positive-weight angle bound in Section~\ref{sec:useful} restricts
to this sector: $\norm{P_E-P_{E_0}}\le\epsilon_n=2/(48n^2-1)$.
Both ranks here are two. Principal-plane rotations give a unitary
$O:E_0\to E$, extended to the sector, with $\norm{O-I}\le2\epsilon_n$.
The two compressed inverses of $\mathcal K'$ differ in norm by at most
$4\epsilon_n\norm{(\mathcal K')^{-1}}\le128\epsilon_n$.
Their inverses have norm at most $K_n$. The inverse-difference identity
therefore yields
\begin{equation}\label{eq:sectorrotation}
 \norm{O^*\mathcal S_E(\mathcal K')O-S_0'}
       \le\frac{256K_n^2}{48n^2-1}.
\end{equation}
Translation symmetry decomposes the full minimization into orthogonal
sectors, so the other characters provide no hidden cross terms.

Now allow every original singleton coefficient, including those outside
$W$. The floor implies $\ker\widehat M'\subset W^\perp$ and
$E\subset\ran\widehat M'$. Hence
\[
 \Sigma'=(P_E(\widehat M')^\dagger P_E|_E)^{-1}
\]
is positive definite. It equals the attained minimum of
$\ip\alpha{\widehat M'\alpha}$ for $P_E\alpha=y$.
Indeed, with $\lambda=\Sigma'y$, the vector
$\alpha=(\widehat M')^\dagger\lambda$ is feasible and every feasible
perturbation is orthogonal to $\lambda$; completing the square proves
minimality. The inequality $\widehat M'\ge\mathcal K'$ gives
$\Sigma'\ge\mathcal S_E(\mathcal K')$, and choosing $\alpha=y$ gives
$\Sigma'\le K_nI$. Combining~\eqref{eq:sectortrace} and
\eqref{eq:sectorrotation}, and dividing the trace by two, proves
\begin{equation}\label{eq:optimizedendpoint}
           \norm{\Sigma'|_{E\cap\mathcal H_1}}\ge b_{\rm end}(n).
\end{equation}
This optimizes all complementary singleton coefficients, rather than
equating a large direct response with a large optimized response.

\subsection{A real field invisible to every endpoint cell}
Suppose $b_{\rm end}(n)>0$. Choose a unit top eigenvector $y$ of the
sector response, with eigenvalue $\beta\ge b_{\rm end}(n)$.
Let $J=n^{5/2}T_L$, $l=J^*y$, $M'=TR'T^*$, and
\[
                       f=-R'T^*(M')^\dagger T_Ll.
\]
The range inclusion just proved and square completion imply
$T(l+f)=0$, $\norm l=1$ and $\ip f{Bf}=c/\beta$.
Both fields lie in $\ran\Xi$ and the actual domain $U$, in sector one.
Since $B\ge13c/3$, $\norm f^2\le3/(13\beta)$. Spectral orthogonality
and~\eqref{eq:lowestenergy} give
\begin{equation}\label{eq:endpointenergy}
 Q(l+f)\le\frac{c\pi^2}{2n^2}+\frac{16c}{13\beta},
 \qquad \norm{l+f}^2\ge1.
\end{equation}
The construction so far may be complex. Sectors $1$ and $N-1$ are
distinct orthogonal characters. All operators are real and preserve
their conjugation pairing. Consequently $\sqrt2\Rea(l+f)$ has exactly
the same norm and energy, and its real low part has norm one.
This gives a nonzero real field without changing any zero observation.

On each first-coordinate slice its point cell is zero by $T(l+f)=0$.
Both full coordinate line sums vanish by $\Xi$. Subtracting the point
from either full line makes each point-by-long cell zero. The whole plane
has sum zero as well, so subtracting the two lines and adding back the
point makes the long-by-long cell zero. The original row normalizations
only multiply these zeros. Thus $C_{\rm end}(l+f)=0$, including every
spatial component. This proves the upper bound in Theorem~\ref{thm:endpoint}.

The lower bound is Corollary~\ref{cor:kernel}, since $f(i)=i$ is a
permutation. The direct sum of channels has the same lower bound and
admits the one-channel upper witness by setting the other channel to zero.
Finally $h_N/\log n\to1$, $K_n/\log n\to48/13$ and
$b_{\rm end}(n)/\log n\to1/(16\pi^2)$. Hence
\[
 \liminf_{n\to\infty}\frac{\lambda_{\rm end}(n)\log n}{c}
 \ge\frac{13}{48},\qquad
 \limsup_{n\to\infty}\frac{\lambda_{\rm end}(n)\log n}{c}
 \le\frac{256\pi^2}{13},
\]
along even sizes. These bounds prove the order, not an exact leading
constant. The sufficient thresholds are extremely conservative and do
not assert a practical small-size counterexample. Any finite penalty on
these same observations vanishes on the constructed field; it cannot
restore the target $c$ at all sufficiently large sizes.

\section{The remaining rows and their exact conditional cost}\label{sec:residual}
For an arbitrary permutation cut, write the original map as $C=(T,P)$,
where $P$ contains all remaining original normalized rows. This notation
does not change their meanings or delete dependent rows. In this section
$P_F$ and $P_L$ denote restrictions of this observation map, not spectral
projections. Put
\[
 M_T=T_FB^{-1}T_F^*,\quad
 f_T(l)=-B^{-1}T_F^*M_T^\dagger T_Ll,\quad
 R_T=B^{-1}-B^{-1}T_F^*M_T^\dagger T_FB^{-1}.
\]
Then define $K_P=P_FR_TP_F^*$ and $r(l)=P_Ll+P_Ff_T(l)$.

\begin{proposition}[Paid residual cost]\label{prop:residual}
For every $l\in L$, $r(l)\in\ran K_P$, and the attained full cancellation
cost equals
\begin{equation}\label{eq:residualcost}
 \min_{C_Ff=-C_Ll}\ip f{Bf}
 =\ip{T_Ll}{M_T^\dagger T_Ll}
       +\ip{r(l)}{K_P^\dagger r(l)}.
\end{equation}
Consequently its comparison with $\ip l{(cI-A_L)l}$ is exactly the
full-kernel target-$c$ criterion.
\end{proposition}
\begin{proof}
Set $X=T_FB^{-1/2}$ and $\Pi_T=I-X^*(XX^*)^\dagger X$.
The finite spectral inverse shows that $\Pi_T$ is the orthogonal projector
onto $\ker X$, and $R_T=B^{-1/2}\Pi_TB^{-1/2}$.
The range theorem provides a full cancelling $f$. The difference
$z=f-f_T(l)$ has $T_Fz=0$, $P_Fz=-r(l)$. With $Y=P_FB^{-1/2}$,
this proves $r(l)\in\ran(Y\Pi_T)=\ran(Y\Pi_TY^*)=\ran K_P$;
the second equality follows by comparing orthogonal kernels.
The minimizer of this residual problem is
$z_*=-R_TP_F^*K_P^\dagger r(l)$, with cost
$\ip{r(l)}{K_P^\dagger r(l)}$, since $R_TBR_T=R_T$.
Whitening by $B^{1/2}$ and completing the square proves minimality.
Also $\ip{f_T(l)}{Bz}=0$ for every $z\in\ker T_F$, as
$Bf_T(l)=-T_F^*M_T^\dagger T_Ll$. The costs therefore add, proving
\eqref{eq:residualcost}. Theorem~\ref{thm:cost} gives the last assertion.
\end{proof}

On $\ran\Xi$ from the preceding section, all long third-coordinate rows
are fixed multiples of their point-plane partners: for any original
second-coordinate interval $I$,
\[
 (Cw)_{I,[1,n-1]}=-\frac{(Cw)_{I,\{0\}}}{\sqrt{n-1}}.
\]
If also $Tw=0$, the prefix and suffix sums on the plane $x_3=0$ are
negatives, since the full second-coordinate line sum is zero. Thus full
invisibility is equivalent to vanishing of the actual planar prefix
sums $\sum_{j<\tau_i}w(i,j,0)$. At endpoint cuts these impose no extra
condition. For moving cuts they can be substantive. Neither their
conditional cost nor a suitable rank deficiency is proved here for
$\tau_i=2i^2\pmod n$.

\section{Limits and the remaining optimization question}\label{sec:limits}
The endpoint theorem proves a matching logarithmic order for the full
original observation kernel, not merely a large dual compression. It
therefore closes the fixed-cut question left open in version~1. The
quadratic moving cut and the optimization over canonical partitions remain
open: their genuine planar prefixes may raise the conditional cost in
\eqref{eq:residualcost}. A generic large direct response cannot decide that
cost, and nominal row counts do not establish a rank defect.

The results concern a finite counting-norm family with prescribed source
weights. No finite-size computation is a proof premise, no formal
verification is claimed, and no continuum limit or physical Yang--Mills
mass gap follows. The finite thresholds are conservative. Establishing
external originality or broad importance would require a more extensive
comparison than the bounded primary-source context given here.
\section*{Model assessments and review context}
At preparation of this revision on 12 November 2026, the
antecedent research sources underwent separate involved analytical source
audits in the same Codex/OpenAI programme context. Those sessions shared
programme history; no provider independence, blind review or human peer
review is claimed. The exact backend model identifiers were unavailable,
and no more specific identities are inferred. The source audits concern
the finite analytical statements, not an exact PDF of this manuscript.

For this new revision at that preparation stage, separate review of its consolidated source,
remote reproducible production, visual PDF inspection and an exact-PDF
scientific assessment had not been completed. Any assessment supplied in
deposit metadata is a separate hash-bound artifact with its own actual
model identity, context and objections. No score, stars, acceptance
decision or assessment of an earlier artifact is transferred to this
version. The full original reports, objections, corrections and context
belong in the accompanying evidence package. AI-assisted preparation and
involved source checks are not represented as independent peer review or
publication acceptance. Lluis Eriksson, the listed author and programme
owner, is also the owner, creator and administrator of AIRR, the intended
deposit destination. This author--platform relationship is not independent
editorial or human review. No affiliation or completed personal review is
invented.

\appendix
\section{Source and build provenance}\label{app:provenance}
This revision consolidates the accepted analytical spectrum, original-row
range and cost, whole-response and useful-response records underlying
version~1, together with the optimized response, conditional cost,
line-average removal and sharp endpoint results. The accompanying
\texttt{source-map289A.json} binds the statements to exact raw source
bytes and records omissions. We include the complete proof chain needed
for the endpoint theorem. The broader dimension-density and general-cut
negative-subspace corollaries of the research notes are omitted, rather
than asserted without their longer proofs. The exact conditional-cost
identity is retained to state precisely what the moving-cut problem needs.
No previously published result or earlier assessment is relabelled as new.
The build system injects the exact frozen source commit before inputting
this file. Its value for this rendering is
\begin{center}\texttt{\SourceCommit}\end{center}
Build logs, reproducibility checks, PDF page-inspection reports and any
exact-PDF assessments are separate hash-bound artifacts. This provenance
statement does not substitute for those records or claim that a particular
review or publication decision has occurred.

\begin{thebibliography}{9}
\bibitem{Eriksson274}
L. Eriksson, \emph{Coverage and coexactness in checkerboard slice observations},
programme working paper, AIRR, version~1 (2026).
\url{https://airr.science/papers/ARR-2026-451FW8N5MC8KYV5S/versions/v1/}.

\bibitem{Eriksson284}
L. Eriksson, \emph{Spectral Reduction and Logarithmic Observation Response
in a Discrete Coexact Model}, AIRR working paper, version~1 (2026).
\url{https://airr.science/papers/ARR-2026-0MKJW2Y51J8AQ96A/versions/v1/}.

\bibitem{Corach}
G. Corach, A. Maestripieri and D. Stojanoff,
\emph{Oblique projections and Schur complements}, arXiv:math/0006120 (2000).
\url{https://arxiv.org/abs/math/0006120}.
\bibitem{Kenyon}
R. Kenyon, \emph{The Laplacian and $\bar\partial$ operators on critical planar
graphs}, arXiv:math-ph/0202018 (2002), especially Theorem~7.3.
\url{https://arxiv.org/abs/math-ph/0202018}.

\bibitem{LawlerLimic}
G. F. Lawler and V. Limic, \emph{Random Walk: A Modern Introduction},
Cambridge University Press, 2010. Author-hosted text, Sections~1.2 and~4.4,
especially Theorem~4.4.4.
\url{https://www.math.uchicago.edu/~lawler/srwbook.pdf}.
\end{thebibliography}
\end{document}
